% Faithful formal-publication transcription: complete Section3.
\section*{3 Communication protocols from tree-like Res(PCd) proofs}
Let $\phi$ be an unsatisfiable CNF formula with $n$ variables. Let us define the communication problem $\oplus_k\operatorname{Search}(\phi)$ with $k$ parties as follows. Assume that the $i$th party has an assignment $\alpha_i\in\{0,1\}^n$ written on the forehead. They aim to find a clause of $\phi$ falsified by the assignment $\sum_{i=1}^k\alpha_i$ (all sums of boolean vectors are computed modulo 2).
\begin{authorclaim}{Lemma1}
Let $\phi$ be an unsatisfiable CNF formula. If there exists a tree-like $\mathrm{Res}(\mathrm{PC}_d)$ proof of $\phi$ of length $m$, then $R^{1/3}_{\mathrm{pub}}(\oplus_{d+1}\operatorname{Search}(\phi))=O(d\cdot\log m)$.
\end{authorclaim}
A slightly weaker version of the following lemma was implicitly proved in [15]:
\begin{authorclaim}{Lemma12 (see proof of Theorem3.11 from [15])}
Let $T$ be a binary tree with $m$ vertices such that the $i$th vertex is labeled with $a_i\in\{0,1\}$ with the hereditary property: for each inner vertex $i$ with direct descendants $c_1$ and $c_2$, if $a_i=1$, then $a_{c_1}=1$ or $a_{c_2}=1$. We also assume that if $r$ is the root of $T$, then $a_r=1$. Assume that we have a one-sided bounded error oracle access to $a_i$ i.e. if we request a value of $a_i$ and $a_i=0$ we get 1 with probability at most $1/2$ and0 with probability at least $1/2$; if $a_i=1$ we get 1 with probability 1. Then there exists an algorithm $A$ that with probability at least $2/3$ returns a leaf $\ell$ of $T$ with $a_\ell=1$ and makes $O(\log m)$ oracle queries to $a_1,\ldots,a_m$.
\end{authorclaim}
Proof. See AppendixA.
\begin{proof}[Proof of Lemma1]
Let $F_1,\ldots,F_m$ be a tree-like $\mathrm{Res}(\mathrm{PC}_d)$-refutation of $\phi$ with the underlying tree $T$, where vertices of $T$ are identified with $[m]$. Then the leaves of $T$ correspond to the clauses of $\phi$ and $m$ is the root of $T$.
Let $\alpha_1,\ldots,\alpha_{d+1}$ be the assignments written on the foreheads of $d+1$ parties. Let $\alpha=\sum_{i=1}^{d+1}\alpha_i$. Let $a_i=1$ iff $\alpha$ falsifies $F_i$ for $i\in[m]$. Then $a_m=1$ since $F_m$ is identically false. For any inner node $v$ of $T$, if $a_v=1$ then for the direct descendants of $v$, $c_1$ and $c_2$ either $a_{c_1}=1$ or $a_{c_2}=1$. In the next paragraphs we show that for any $i\in[m]$ there exists a NOF $(d+1)$-party protocol that computes $a_i$ given that for each $j\in[d+1]$ the $j$th party has $\alpha_j$ written on their forehead such that
\begin{itemize}
\item the protocol transmits $O(d)$ bits;
\item the protocol has one-sided bounded error: if $a_i=1$ then the protocol returns 1 with probability 1 and if $a_i=0$ the protocol returns 0 with probability at least $1/2$.
\end{itemize}
Then we use this protocol to compute $a_i$ as an oracle in the algorithm given by Lemma12 and thus show that there is a NOF $(d+1)$-party protocol computing $\oplus_{d+1}\operatorname{Search}(\phi)$ with communication cost $O(d\log m)$.
Now we show that for every $\ell\in[m]$, $F_\ell(\alpha)$ can be computed by a $(d+1)$-party NOF protocol with one-sided error using $O(d)$ bits of communication. Let $F_\ell=\bigvee_{j=1}^t(f_j=1)$, where $f_1,\ldots,f_t$ are polynomials over $\mathbb F_2$ of degree at most $d$. Let $z_1,\ldots,z_n$ be the variables of $\phi$. Let us introduce new variables $y_{1,1},\ldots,y_{1,n},\ldots,y_{d+1,1},\ldots,y_{d+1,n}$ and assume that for each $i\in[d+1]$ the $i$th party has the value of variables $y_{i,1},y_{i,2},\ldots,y_{i,n}$ written on the forehead or in other words $\alpha_i$ assigns values of $y_{i,1},y_{i,2},\ldots,y_{i,n}$. Let $\bar f_j$ denote $f_j$ after substitution $z_\ell:=y_{1,\ell}+y_{2,\ell}+\cdots+y_{d+1,\ell}$ for $\ell\in[n]$; $j\in[t]$. Since for all $j\in[t]$, $\deg\bar f_j=\deg f_j\le d$, we can represent $\bar f_j=\bar f_j^{(1)}+\cdots+\bar f_j^{(d+1)}$ such that $\bar f_j^{(s)}$ does not contain variables $y_{s,1},\ldots,y_{s,n}$ for each $s\in[d+1]$. Then the $i$th party can compute $\bar f_1^{(i)}(\alpha_1,\ldots,\alpha_{d+1}),\ldots,\bar f_t^{(i)}(\alpha_1,\ldots,\alpha_{d+1})$. Notice that $F_\ell=\neg(\bigwedge_{j=1}^t(f_j=0))$.
The final step of the protocol exploits the idea used to construct a short randomized communication protocol for equality. Take a random uniformly distributed vector $(e_1,\ldots,e_t)\in\mathbb F_2^t$. Then all parties compute
\[\sum_{j=1}^te_jf_j(\alpha)=\sum_{i=1}^{d+1}\underbrace{\sum_{j=1}^te_j\bar f_j^{(i)}}_{i\text{th party}}\]
with $O(d)$ bits of communication and the protocol halts. To bound the error probability we use the following well-known statement:
\begin{authorclaim}{Proposition13 (Random subsum principle)}
For any $x\in\mathbb F_2^k\setminus\{0^k\}$,
\[\Pr_{y\leftarrow U(\mathbb F_2^k)}\left[\sum_{i=1}^ky_ix_i=1\right]=\frac12.\]
\end{authorclaim}
If $F_\ell(\alpha)=1$ then $\Pr[\sum_{j=1}^te_jf_j(\alpha)\ne0]=1/2$ by the random subsum principle. If $F_\ell(\alpha)=0$, then $\Pr[\sum_{j=1}^te_jf_j(\alpha)=0]=1$.
\end{proof}
\begin{authorclaim}{Remark14}
Similarly to the proof of Lemma1 one can prove that if an unsatisfiable CNF formula $\phi$ has a tree-like $\mathrm T^{\mathrm{cc}}_{\mathrm{os}}(k,c)$ refutation of length $\ell$, then for any $k$-partition of the variables, there is a randomized bounded-error $k$-party NOF protocol for Search$(\phi)$ with communication cost $O(c\log\ell)$. Thus, the bound from Lemma9 can be slightly improved in the case of one-sided error.
\end{authorclaim}
